\section{Conclusions}
\label{iii:sec:conclusiones}
\begingroup
\fontsize{10.5}{12.8}\selectfont
\setlength{\parskip}{2pt plus .5pt minus .5pt}

Le résultat central est une détermination constitutive réversible dans
le domaine contractant : les deux canaux orientés d’un transport
produisent quatre coordonnées reliées et ces coordonnées permettent de
retrouver la paire angulaire. La structure fixe d’abord l’équation de
réponse ; son évaluation détermine les valeurs sans perdre l’orientation
nécessaire à la reconstruction de leurs antécédents spectraux. Cette double
direction, constructive et inverse, donne un contenu précis à la thèse
structure, équation et valeur.

\subsection*{La réponse conserve la paire angulaire}

Les incidences $90$ et $120$ déterminent une unique solution de
l’équation de complétion pour le transport fixé. Sa fonction scalaire
est une bijection strictement décroissante de $(0,1)$ sur $(3/4,1)$.
Par conséquent, l’équation cubique~\eqref{iii:eq:cubic} possède exactement
une racine admissible pour chaque réponse de canal. L’inversion n’est pas
une coïncidence numérique : elle reconstruit les deux contractions et, par
leurs logarithmes, la paire angulaire.

Les quatre coordonnées du vide satisfont
\[
\widehat\mu\widehat\varepsilon=\widehat c^{-2},\qquad
\widehat\mu/\widehat\varepsilon=\widehat Z^2.
\]
Dans les règles fixées de produit, de quotient et de positivité, ces
identités déterminent les évaluations quadratiques. Leur dépendance
commune se manifeste également sous l’échange de feuillets : la permittivité et
la perméabilité se permutent, l’impédance s’inverse et la vitesse est
conservée. Le produit symétrique et le quotient orienté transmettent ainsi
des informations complémentaires. La récupération maintient les étiquettes
de feuillet et la normalisation interne ; l’état d’origine conserve
en outre le chemin, la frontière et la mémoire.

La factorisation $\mathcal V=a\exp(\eta\mathcal R)$ identifie
ces deux informations dans l’opérateur lui-même. Le facteur positif
$a=\sqrt{\det\mathcal V}$ détermine la vitesse normalisée ;
la composante unimodulaire détermine l’impédance. Conserver les deux
permet de reconstruire la réponse. Cette séparation explique la relation
entre les feuillets réciproques de produit unitaire et les quatre lectures
du vide : la normalisation sépare le mode commun, qui demeure
nécessaire pour retrouver l’opérateur complet.

L’inverse permet de relier des réponses qui évaluent des fonctions distinctes
des mêmes canaux. Le rapport de leurs exposants reconstruit
l’anisotropie de l’ellipse d’action et l’impédance LC relative,
selon la proposition~\ref{iii:prop:forma-LC}. La section d’action
et les références de fréquence et d’impédance complètent leurs grandeurs dimensionnelles.
L’identité $c_{\rm int}=c_{\rm vac}$ établit, à son tour, que le
lecteur cinématique et le constitutif réalisent la même section de
vitesse dans une coordinatisation commune. La trace normalisée conserve cette
lecture sous amplification nonadique.

\subsection*{Composition et reconstruction entre secteurs}

La loi transportée $r\odot t=f(\psi(r)\psi(t))$ exprime la composition
du transport en coordonnées constitutives. Son élément neutre est $3/4$ dans
l’extension de frontière, et $-\log\psi(r)/30$ transforme cette composition en
somme. Les projecteurs communs préservent la réalisation opératorielle et les
étiquettes de feuillet. La vitesse d’un état unique conserve, pour sa part,
le produit ordinaire de ses deux canaux. Les deux opérations demeurent
distinguées par leurs arguments et par la preuve exacte de leur relation.

La reconstruction Higgs--impédance établit une autre connexion explicite.
Les couplages de jauge restituent la contraction commune et le
couplage électromagnétique ; l’auto-couplage de Higgs restitue la
coordonnée orientée. Dans le domaine de la réalisation à l’arbre, ces
données déterminent les deux canaux et leurs réponses du vide. Les
inégalités de domaine et les compatibilités généalogiques précisent
quels triplets appartiennent à la réalisation considérée. Conserver les
projecteurs marqués permet de restituer également l’opérateur ; l’histoire
enrichie maintient ses données de provenance supplémentaires.

\subsection*{Le cycle relie réponse, moment et orientation}

La réalisation dodécaphasique explique la provenance de la fonction
rationnelle par un quotient de déterminants en dimensions $3$ et $4$.
La différence des rangs produit $-1/12$ comme trace normalisée dans ce
domaine fini. Le plan du quart de tour conserve une autre information :
son opérateur satisfait $J^2=-I$ et son coefficient de résolvante retient
le signe d’orientation. Les deux intégrations logarithmiques de ce
coefficient, avec leurs conditions initiales, déterminent le moment
de Catalan à profondeur illimitée.

La table des traces $b_n$ reconstruit $-\log f(s)$ et produit des moments
convergents pour $\operatorname{Re}z>0$, avec une borne explicite du
reste. Sa valeur à la profondeur un est $\log(4/3)$, compatible avec
l’extrémité constitutive. La représentation de Mellin et la distinction
entre produit diagonal et tensoriel conservent les opérations qui
précèdent l’évaluation scalaire de chaque moment.

La restitution intégrale du
théorème~\ref{iii:thm:restitucion-funcional-catalan} clôt la
correspondance fonctionnelle dans les deux sens :
\[
 \mathcal C_2(s)=\int_0^s\log\frac{s}{t}\,
       \frac{1+t}{1+t+t^2}f(t)\,dt,\qquad
 f(s)=\frac{1+s+s^2}{s(1+s)}\mathcal D^2\mathcal C_2(s).
\]
La condition $\mathcal C_2(s)\to0$ à l’origine élimine toute
liberté additive et logarithmique. La valeur extrême $G_{\rm Cat}$,
le noyau orienté et la réponse constitutive se trouvent ainsi situés
dans une même collection, avec des opérations de restitution explicites.
Les arguments $s_\pm$ conservent leur détermination angulaire ; la trace
normalisée maintient la vitesse du même état en la reliant
à l’action et à la longueur. Cette correspondance spécifie le
contenu structurel qui est conservé lors du passage de la collection à ses valeurs.

L’identité Catalan--Pell précise une évaluation algébrique de cette
collection :
\[
\mathcal C_2(2-\sqrt3)
=\frac23G_{\rm Cat}-\frac{\pi}{12}\log(2+\sqrt3).
\]
Le coefficient $2/3$ provient des classes impaires du calendrier et
du poids de profondeur ; les deux premières dérivées logarithmiques dans cette section prennent les valeurs
$\pi/12$ et $1/4$. La réciprocité de l’unité de Pell relie
expansion et contraction, tandis que les deux arguments constitutifs
restent déterminés par la paire angulaire. Cette relation maintient
visibles la collection, ses lecteurs et les sections où ils sont évalués.

\Needspace{9\baselineskip}
La fonctionnelle de Barbero--Immirzi s’exprime ensuite comme différence
orientée d’une même fonction des canaux reconstruits :
\[
\gamma=\mathcal H(\psi(r_-))-\mathcal H(\psi(r_+)).
\]
L’incidence et le retour dodécaphasique fixent ses termes ;
l’inversion constitutive permet de les transporter vers la réponse du
vide. La marque d’orientation étant fixée, l’échange des canaux
change le signe de $\gamma$ ; la conjugaison conjointe conserve la
fonctionnelle. Cette distinction prouve que son contenu dépasse le spectre
sans étiquettes et se maintient sous les amplifications normalisées.

% BEGIN SINTESIS_ADITIVA_20260922
\par
La section~\ref{delta22:iii:restitucion} réunit le potentiel spectral de la réponse constitutive, la restitution globale des canaux au moyen de la vitesse normalisée et de Barbero--Immirzi, ainsi que ses bornes de stabilité. Sur l’image engendrée, avec les étiquettes et les bases conservées, la composition retrouve la paire angulaire, la section d’action et la lecture périodique du registre dodécaphasique.
% END SINTESIS_ADITIVA_20260922

La résolution harmonique du motif dodécaphasique conserve exactement
les modes $3,9,4,8$, avec les coefficients non normalisés $72,72,-6,-6$.
Le rapport des modules $12:1$ exprime les multiplicités orientées.
Le défaut d’incidence $-1/12$ et le coefficient singulier $1/24$
restent liés à leurs propres opérations. La comparaison explicite
avec le lecteur à poids égaux, dont le coefficient est $1/1080$,
précise comment la chaîne orientée intervient dans la normalisation
du retour sans remplacer les fonctionnelles par leurs pôles.

\subsection*{Du bilan arithmétique aux relations électriques}

La discrépance des vacances fournit une polarisation centrée dont la
grandeur reste inférieure à une unité et dont l’accroissement satisfait
un bilan temporel exact. Les changements entre les conventions de
plancher et de plafond sont enregistrés par un terme de frontière. Une
section de charge et une durée orientée réalisent ce bilan comme
charge centrée et courant, en conservant sa loi télescopique.

Le préfacteur électronique conserve également sa structure antécédente.
L’énumération des fibres dans $D_9^3$ fixe le mode singulier tritique
de valeur $1/2$ ; le spectre complet de son appariement prouve
\[
\|[P_{\Sigma,3},P_{\Pi,3}]\|=\sqrt3/4.
\]
Son défaut quadratique $3/4$ coïncide avec la limite du quotient des
secteurs $90/120$. Cette égalité relie des résultats d’opérateurs
spécifiés, chacun avec son espace et son action : le préfacteur est
une norme d’incompatibilité arithmétique, tandis que la réponse
intérieure dépend de ses canaux contractifs.

La section positive de charge déterminée par
$Z_{\rm vac}e_a^2=2\alpha h_a$ réunit le couplage, l’action et
l’impédance. Les relations de von Klitzing, de conductance, de flux
et de Josephson satisfont exactement
\[
R_KG_0=2,\qquad G_0Z_{\rm vac}=4\alpha,\qquad
K_{J,a}\Phi_{0,a}=1,\qquad K_{J,a}^2R_K=4/h_a.
\]
Le retour de l’action conserve résistance et conductance ; charge et
flux se transforment avec $R_{\rm act}^{1/2}$ et Josephson avec
$R_{\rm act}^{-1/2}$. Ces dépendances expliquent conjointement
quelles grandeurs retiennent l’action et lesquelles conservent uniquement leur
quotient constitutif.

La réalisation dimensionnelle préserve toutes ces identités lorsque
les bases d’action, de charge et de vitesse changent. Les valeurs
adimensionnelles obtenues appartiennent à la mise en coordonnées interne de la
construction ; leur reconnaissance métrologique exige l’évaluation
indépendante des sections dimensionnelles et du régime physique
correspondant. La portée mathématique établie est la réponse
positive et inversible, ses compositions et sa covariance dans les
domaines déclarés.

L’ensemble montre comment une même généalogie APP--TRIT--TPK conserve
les relations nécessaires pour passer de chemins et d’opérations
arithmétiques à des canaux, des moments, une réponse constitutive et des relations
électriques. Les valeurs acquièrent ainsi une provenance expliquée et
une information récupérable : elles expriment les opérations qui les
produisent et permettent de reconnaître, au sein de leur représentation, la
structure qui demeure en elles.
\par\endgroup
